\section{What counts as progress? An operational definition}
\label{sec:definition}

We need a definition of stagnation that (i) a human can apply to a trajectory, (ii) a runtime can
approximate without the reference solution, and (iii) is falsifiable, in the sense that data could show
the definition is not worth monitoring. The definition below is the one handed to annotators before any
monitor was written; the codebook is reproduced in Appendix~\ref{app:codebook}.

\subsection{Windows, not steps}
A single step is rarely informative: reading one file can be the most productive step of a run, running
one test can be wasted motion. We therefore judge \emph{windows} of $w$ consecutive steps
($w=\pwParam$ throughout) and ask, at the end of each window, whether the trajectory moved.

\subsection{Three channels of progress}
A window has advanced the task if at least one of the following is true, judged from the trajectory
itself:

\begin{description}[leftmargin=1.1em,style=nextline,itemsep=2pt]
\item[E --- epistemic.] Something new \emph{and relevant to this task} became known: a file, function,
symbol, error signature or test name plausibly in scope of the task statement; a hypothesis eliminated
by evidence; the search space narrowed. Irrelevance is judged against the task statement, the reported
failure, and the artefact the task asks for.
\item[I --- implementation.] A plausible artefact changed durably --- the change survives the window
rather than being immediately reverted --- or a previously malformed artefact became well-formed.
\item[V --- verification.] Objective evidence about correctness improved: a failing test passes, the
failure count falls, a build moves from failing to succeeding, or the error signature advances past the
previous blocker.
\end{description}

\noindent A window in which none of the three advanced \emph{while the agent kept expending actions} is
\emph{stagnant}. Four additional labels keep the binary target from swallowing distinct phenomena:
\textsc{regression} (a previous success was undone), \textsc{done\_redundant} (the task was already
complete and verified and the agent keeps working without a new goal), \textsc{blocked\_external}
(progress is impossible for reasons outside the agent's control, e.g.\ a dependency download in
flight), and \textsc{uncertain}. For the binary task, \textsc{stagnant} and \textsc{done\_redundant} are
positive, \textsc{productive} and \textsc{regression} are negative, and the other two are held out and
counted.

\subsection{Why this definition, and what would falsify it}
The definition is deliberately conservative in one direction and generous in another. It never calls a
window stagnant merely because commands repeat: an edit loop that finally produces passing tests is
progress, and a repeated test invocation whose failure output shrinks is progress. It counts purely
epistemic work --- looking, understanding, narrowing --- as progress even with an untouched workspace,
which is exactly the case that workspace-change monitors miss. It would be falsified as a research target
if the channels it names were unobservable at runtime, or if the labels it produces had no relationship to
task outcomes or to cost; Section~\ref{sec:results} reports both checks.

The annotator sees the whole trajectory, because deciding whether information was already known or whether
a change survived requires the future. The monitor, at step $t$, sees only the task statement, the actions
and observations of steps $\le t$, and the diffs implied by those actions: never the reference patch, the
hidden tests, later steps, the final reward, or any annotation attached to a later window. Every feature
used in Section~\ref{sec:method} is computed from a window ending at $t$; corpus-level normalisation
statistics are estimated on the training split only, and the split is by task, so no monitor can memorise a
task.
